\originalchapter{历年第一次期中考试}
\label{chap:exams-midterm1}
本章收录官方考试目录中的九份第一次期中考试, 依原年度与季节排列. 每卷均为50分钟、100分. 原考试要求清楚展示计算, 使用已知事实时说明事实及其适用理由, 不使用计算器、书籍或笔记, 尽量化简答案; 阶乘不必展开, 2018--2019年的试卷也明确允许保留二项式系数. 以下解答均\textbf{依据官方解答重构, 补齐推导}. 原文的数学误印在相应解答中明示; 题面重复字母保留, 不改作新的题号.

\originalsection{2011年春季}
% SOURCE 2011-spring-midterm1
\noindent 官方来源: \href{https://ocw.mit.edu/courses/18-600-probability-and-random-variables-fall-2019/93a853eac855b4de3cbee2b7fa1a761e_MIT18_600F19_mid1_S2011.pdf}{原试卷PDF}; \href{https://ocw.mit.edu/courses/18-600-probability-and-random-variables-fall-2019/61394f8d49ead2d89ee46e87b07a5a14_MIT18_600F19_mid1_S2011_soln.pdf}{官方解答PDF}.
\begin{exercise}\label{ex:m1-2011s-q1}
\textbf{2011年春季, 原题1 (20分).} 无限次独立投掷一枚硬币, 每次正面概率为$p$.
\begin{enumerate}
\item[(a)]\phantomsection\label{ex:m1-2011s-q1-a} 设首次正面出现在第$X$次投掷, 即$X$为得到一次正面所需的投掷次数. 用$p$表示$\E[X]$.
\item[(b)]\phantomsection\label{ex:m1-2011s-q1-b} 用$p$表示前10次中恰有5次正面的概率.
\item[(c)]\phantomsection\label{ex:m1-2011s-q1-c} 用$p$表示第5次正面恰出现在第10次投掷的概率.
\end{enumerate}
\end{exercise}
\begin{solution}
\textbf{(a)}\phantomsection\label{sol:m1-2011s-q1-a} 当$0<p\le1$时, $\{X>n\}$表示前$n$次全是反面, 概率为$(1-p)^n$. 用正整数值随机变量的尾和公式,
\[
\E[X]=\sum_{n=0}^{\infty}\Prob(X>n)=\sum_{n=0}^{\infty}(1-p)^n=\frac1p.
\]
这里尾和来自逐点恒等式$X=\sum_{n\ge0}\ind_{\{X>n\}}$, 非负项可逐项取期望. 若$p=0$, 永无正面, 有限等待时间不存在; 将$X$取为$+\infty$时, 期望也是$+\infty$.

\textbf{(b)}\phantomsection\label{sol:m1-2011s-q1-b} 在10个位置中选出正面的5个位置, 每种指定序列的概率是$p^5(1-p)^5$. 互斥序列求和得到$\binom{10}{5}p^5(1-p)^5$.

\textbf{(c)}\phantomsection\label{sol:m1-2011s-q1-c} 第10次必须为正面, 前9次必须恰有4个正面. 前9次与第10次独立, 所以所求为$\binom94p^4(1-p)^5p=\binom94p^5(1-p)^5$. 后两式在$p=0,1$时也给出正确的零概率.
\end{solution}

\begin{exercise}\label{ex:m1-2011s-q2}
\textbf{2011年春季, 原题2 (20分).} Jill向在monster.com上找到的1000家公司发送简历. 每家公司以$3/1000$的概率回复, 各公司的行为相互独立. 设$R$为回复的公司数.
\begin{enumerate}
\item[(a)]\phantomsection\label{ex:m1-2011s-q2-a} 计算$\E[R]$, 给出精确值而非近似值.
\item[(b)]\phantomsection\label{ex:m1-2011s-q2-b} 计算$\Var[R]$, 给出精确值而非近似值.
\item[(c)]\phantomsection\label{ex:m1-2011s-q2-c} 用泊松随机变量近似估计$\Prob(R=3)$.
\end{enumerate}
\end{exercise}
\begin{solution}
\textbf{(a)}\phantomsection\label{sol:m1-2011s-q2-a} 令$I_j$表示第$j$家公司是否回复, 则$R=\sum_{j=1}^{1000}I_j$, $\E[I_j]=3/1000$. 期望可加, 故$\E[R]=1000(3/1000)=3$.

\textbf{(b)}\phantomsection\label{sol:m1-2011s-q2-b} $I_j^2=I_j$, 因而$\Var(I_j)=p-p^2=p(1-p)$, $p=3/1000$. 各指标独立使协方差为零, 所以$\Var(R)=1000p(1-p)=3(1-3/1000)=2991/1000$.

\textbf{(c)}\phantomsection\label{sol:m1-2011s-q2-c} 精确分布为$\Bin(1000,3/1000)$. 次数多且单次回复率小, 用均值相同的$\Pois(3)$近似, 得$\Prob(R=3)\approx\ee^{-3}3^3/3!=9\ee^{-3}/2$. 这是近似概率, 前两问仍是精确矩.
\end{solution}

\begin{exercise}\label{ex:m1-2011s-q3}
\textbf{2011年春季, 原题3 (10分).} 非负整数四元组$(x_1,x_2,x_3,x_4)$满足$x_1+x_2+x_3+x_4=10$, 这样的四元组有多少个?
\end{exercise}
\begin{solution}
\phantomsection\label{sol:m1-2011s-q3}
把10颗星排成一行, 插入3条分隔线, 四段星的个数依次为$x_1,x_2,x_3,x_4$. 分隔线相邻或位于端点允许某段为零, 因而与非负整数解一一对应. 在13个位置中选3个放分隔线, 解数为$\binom{13}{3}=286$.
\end{solution}

\begin{exercise}\label{ex:m1-2011s-q4}
\textbf{2011年春季, 原题4 (10分).} 一张彩票以百万分之一的概率使你赢得一百万美元. 设$X$为赢得的金额. 计算:
\begin{enumerate}\item[(a)]\phantomsection\label{ex:m1-2011s-q4-a} $\E[X]$. \item[(b)]\phantomsection\label{ex:m1-2011s-q4-b} $\Var[X]$.\end{enumerate}
\end{exercise}
\begin{solution}
\textbf{(a)}\phantomsection\label{sol:m1-2011s-q4-a} 奖金只取$10^6$和$0$, 所以$\E[X]=10^6\cdot10^{-6}=1$美元. 本题问奖金, 没有给彩票购买价格, 故不扣票价.

\textbf{(b)}\phantomsection\label{sol:m1-2011s-q4-b} 二阶矩$\E[X^2]=(10^6)^2\cdot10^{-6}=10^6$, 因而$\Var(X)=\E[X^2]-\E[X]^2=10^6-1$. 方差的单位为美元平方.
\end{solution}

\begin{exercise}\label{ex:m1-2011s-q5}
\textbf{2011年春季, 原题5 (20分).} 连续随机变量$X$的密度为
\[
f_X(x)=\begin{cases}2x,&0\le x\le1,\\0,&x\notin[0,1].\end{cases}
\]
计算: \begin{enumerate}\item[(a)]\phantomsection\label{ex:m1-2011s-q5-a} $\E[X]$.\item[(b)]\phantomsection\label{ex:m1-2011s-q5-b} $\Var[X]$.\item[(c)]\phantomsection\label{ex:m1-2011s-q5-c} 累积分布函数$F_X$.\end{enumerate}
\end{exercise}
\begin{solution}
密度非负且$\int_0^1 2x\dd x=1$, 因而给出了概率分布.

\textbf{(a)}\phantomsection\label{sol:m1-2011s-q5-a} 密度加权平均给出$\E[X]=\int_0^1x\cdot2x\dd x=2/3$.

\textbf{(b)}\phantomsection\label{sol:m1-2011s-q5-b} 同理$\E[X^2]=\int_0^1x^2\cdot2x\dd x=1/2$, 故$\Var(X)=1/2-(2/3)^2=1/18$.

\textbf{(c)}\phantomsection\label{sol:m1-2011s-q5-c} 由$F_X(a)=\Prob(X\le a)=\int_{-\infty}^a f_X(x)\dd x$, 在支撑区间两侧及其内部依次得到
\[
F_X(a)=\begin{cases}0,&a<0,\\a^2,&0\le a\le1,\\1,&a>1.\end{cases}
\]
在$0,1$处两侧值一致, 符合连续分布没有端点原子的性质.
\end{solution}

\begin{exercise}\label{ex:m1-2011s-q6}
\textbf{2011年春季, 原题6 (20分).} 标准52张牌中有4张A. 随机排列整副牌, 所有$52!$种排列等可能. 计算:
\begin{enumerate}\item[(a)]\phantomsection\label{ex:m1-2011s-q6-a} 最上面4张全为A的概率.\item[(b)]\phantomsection\label{ex:m1-2011s-q6-b} 最上面4张都不是A的概率.\item[(c)]\phantomsection\label{ex:m1-2011s-q6-c} 最上面4张中A的张数的期望. 答案有简单形式.\end{enumerate}
\end{exercise}
\begin{solution}
\textbf{(a)}\phantomsection\label{sol:m1-2011s-q6-a} 前4张由4张A任意排列, 有$4!$种, 其余48张有$48!$种. 故概率为$4!48!/52!=1/\binom{52}{4}$.

\textbf{(b)}\phantomsection\label{sol:m1-2011s-q6-b} 依次选择前4张非A牌的选择数为$48\cdot47\cdot46\cdot45$, 剩余48张任意排列. 概率为
\[
\frac{48\cdot47\cdot46\cdot45\cdot48!}{52!}
=\frac{\binom{48}{4}}{\binom{52}{4}}.
\]
顺序计数和无序取样计数在此给出同一概率.

\textbf{(c)}\phantomsection\label{sol:m1-2011s-q6-c} 令$I_j$为第$j$张是A的指标, $j=1,2,3,4$. 每个位置的牌均匀, 所以$\E[I_j]=4/52=1/13$. A的总数为四个指标之和, 故期望为$4/13$. 指标并不独立, 但期望可加无需独立性.
\end{solution}

\originalsection{2011年秋季}
% SOURCE 2011-fall-midterm1
\noindent 官方来源: \href{https://ocw.mit.edu/courses/18-600-probability-and-random-variables-fall-2019/1419e4cab0601d2723040c9b58427a79_MIT18_600F19_mid1_F2011.pdf}{原试卷PDF}; \href{https://ocw.mit.edu/courses/18-600-probability-and-random-variables-fall-2019/d82587d0909f829fb4b17d4e0b240ea3_MIT18_600F19_mid1_F2011_soln.pdf}{官方解答PDF}.
\begin{exercise}\label{ex:m1-2011f-q1}
\textbf{2011年秋季, 原题1 (20分).} Jill去钓鱼. 她钓鱼的每一分钟, 以$1/600$的概率捕到一条鱼, 各分钟相互独立. 她连续钓鱼15小时, 即900分钟. 设$N$为捕获的鱼数.
\begin{enumerate}\item[(a)]\phantomsection\label{ex:m1-2011f-q1-a} 精确计算$\E[N]$及$\Var[N]$.\item[(b)]\phantomsection\label{ex:m1-2011f-q1-b} 精确计算恰捕到3条鱼的概率.\item[(c)]\phantomsection\label{ex:m1-2011f-q1-c} 用泊松随机变量近似计算恰捕到3条鱼的概率.\end{enumerate}
\end{exercise}
\begin{solution}
\textbf{(a)}\phantomsection\label{sol:m1-2011f-q1-a} 令每分钟的捕鱼指标为$I_j$, 则$N=\sum_{j=1}^{900}I_j\sim\Bin(900,1/600)$. 指标的均值为$p$, 方差为$p(1-p)$, 不同分钟独立, 因此
\[
\E[N]=\frac{900}{600}=\frac32,\qquad
\Var(N)=900\frac1{600}\frac{599}{600}=\frac{599}{400}.
\]
\textbf{(b)}\phantomsection\label{sol:m1-2011f-q1-b} 从900分钟中选3分钟捕到鱼, 其余897分钟没有鱼, 所得互斥序列总概率为$\binom{900}{3}(1/600)^3(599/600)^{897}$.

\textbf{(c)}\phantomsection\label{sol:m1-2011f-q1-c} 用$\Pois(3/2)$近似$N$, 得$\Prob(N=3)\approx\ee^{-3/2}(3/2)^3/3!=9\ee^{-3/2}/16$.
\end{solution}

\begin{exercise}\label{ex:m1-2011f-q2}
\textbf{2011年秋季, 原题2 (10分).} 房间里有10个人. 假设将出生月份分配给这10人的全部$12^{10}$种方式等可能, 求没有任何两人出生于同一个月份的概率.
\end{exercise}
\begin{solution}
\phantomsection\label{sol:m1-2011f-q2}
把10个人固定排序. 第一人的月份有12种选择, 第二人要避免第一人的月份, 有11种, 依此到第10人还有3种. 无碰撞的有序月份序列数为$12\cdot11\cdots3=12!/2!$. 因此概率为$12!/(2!12^{10})=\binom{12}{10}10!/12^{10}$. 此计算使用题给均匀模型, 不涉及现实月份出生频率.
\end{solution}

\begin{exercise}\label{ex:m1-2011f-q3}
\textbf{2011年秋季, 原题3 (10分).} $X,Y,Z$相互独立, 每个都以概率$1/2$取$0$, 以概率$1/2$取$1$.
\begin{enumerate}\item[(a)]\phantomsection\label{ex:m1-2011f-q3-a} 计算$\Prob(X+Y+Z=1\mid X-Y=0)$.\item[(b)]\phantomsection\label{ex:m1-2011f-q3-b} 事件$\{X=Y\},\{Y=Z\},\{X=Z\}$是否相互独立? 是否两两独立? 说明理由.\end{enumerate}
\end{exercise}
\begin{solution}
\textbf{(a)}\phantomsection\label{sol:m1-2011f-q3-a} 八个三元组$(X,Y,Z)$等可能. 条件$X=Y$允许$(0,0,0),(0,0,1),(1,1,0),(1,1,1)$四种; 其中和为1的只有$(0,0,1)$. 因而条件概率为$(1/8)/(4/8)=1/4$.

\textbf{(b)}\phantomsection\label{sol:m1-2011f-q3-b} 记三个事件为$A,B,C$. 各自概率为$1/2$. 任意两个相等条件都迫使$X=Y=Z$, 只剩全零或全一, 故$\Prob(AB)=\Prob(AC)=\Prob(BC)=1/4=(1/2)^2$. 所以两两独立. 但$\Prob(ABC)=1/4\ne(1/2)^3$, 因而不相互独立; 检查二事件交不足以检查三事件独立性.
\end{solution}

\begin{exercise}\label{ex:m1-2011f-q4}
\textbf{2011年秋季, 原题4 (20分).} 无限次独立投掷一枚硬币, 每次正面概率为$p$.
\begin{enumerate}\item[(a)]\phantomsection\label{ex:m1-2011f-q4-a} 首次正面出现在第$X$次, 即得到一次正面所需次数为$X$. 用$p$表示$X$的期望及方差.\item[(b)]\phantomsection\label{ex:m1-2011f-q4-b} 第5次正面出现在第$Y$次. 用$p$表示$Y$的期望及方差.\end{enumerate}
\end{exercise}
\begin{solution}
\textbf{(a)}\phantomsection\label{sol:m1-2011f-q4-a} 设$0<p\le1$, $q=1-p$. 前面习题\ref{ex:m1-2011s-q1}已用尾和得到$\E[X]=1/p$. 为算二阶矩, 利用对正整数$k$成立的$k^2=\sum_{n=0}^{k-1}(2n+1)$, 得
\[
\E[X^2]=\sum_{n=0}^{\infty}(2n+1)q^n
=\frac{2q}{(1-q)^2}+\frac1{1-q}
=\frac{2-p}{p^2}.
\]
其中$\sum nq^n=q/(1-q)^2$由几何级数在$|q|<1$内部求导得到. 于是$\Var(X)=(2-p)/p^2-1/p^2=(1-p)/p^2$.

\textbf{(b)}\phantomsection\label{sol:m1-2011f-q4-b} 将第5次正面前的投掷分成五个等待段, 每段从上一个正面之后到下一个正面为止. 各段长度$G_1,\ldots,G_5$独立, 都有与$X$相同的几何分布: 对任意给定正整数段长, 相应不交的投掷段的联合概率分解为$\prod_{j=1}^5q^{G_j-1}p$. 因而$Y=G_1+\cdots+G_5$, 期望与独立和的方差分别给出$\E[Y]=5/p$, $\Var(Y)=5(1-p)/p^2$. 若$p=1$, 等待次数分别恒为1和5, 方差为零; 若$p=0$, 两个有限等待时间均不存在.
\end{solution}

\begin{exercise}\label{ex:m1-2011f-q5}
\textbf{2011年秋季, 原题5 (20分).} 连续随机变量$X$的密度在$x>0$时为$\ee^{-x}$, 在$x<0$时为$0$. 计算:
\begin{enumerate}\item[(a)]\phantomsection\label{ex:m1-2011f-q5-a} $\E[X]$.\item[(b)]\phantomsection\label{ex:m1-2011f-q5-b} $\Prob(X\in[-50,50])$.\item[(c)]\phantomsection\label{ex:m1-2011f-q5-c} 累积分布函数$F_X$.\end{enumerate}
\end{exercise}
\begin{solution}
密度在单点$0$的取值不影响积分或分布, 原题未指定此值无需补成概率条件.

\textbf{(a)}\phantomsection\label{sol:m1-2011f-q5-a} 分部积分得$\E[X]=\int_0^\infty x\ee^{-x}\dd x=[-x\ee^{-x}]_0^\infty+\int_0^\infty\ee^{-x}\dd x=1$, 因为$x\ee^{-x}\to0$.

\textbf{(b)}\phantomsection\label{sol:m1-2011f-q5-b} 负半轴没有质量, 故$\Prob(-50\le X\le50)=\int_0^{50}\ee^{-x}\dd x=1-\ee^{-50}$.

\textbf{(c)}\phantomsection\label{sol:m1-2011f-q5-c} 对$a\le0$, $F_X(a)=0$; 对$a>0$, 积分到$a$得$F_X(a)=1-\ee^{-a}$. 这两个表达式在$a=0$处相接, 在$a\to\infty$时趋于1.
\end{solution}

\begin{exercise}\label{ex:m1-2011f-q6}
\textbf{2011年秋季, 原题6 (20分).} 52个人标为$1,2,\ldots,52$, 将帽子投入箱中混合后每人随机得到一顶, 全部$52!$种排列等可能. 计算:
\begin{enumerate}\item[(a)]\phantomsection\label{ex:m1-2011f-q6-a} 前26个人都拿到自己帽子的概率.\item[(b)]\phantomsection\label{ex:m1-2011f-q6-b} 存在26对帽子相互交换的人的概率: 每对可标为$(a,b)$, $a$拿$b$的帽子, $b$拿$a$的帽子.\end{enumerate}
\end{exercise}
\begin{solution}
\textbf{(a)}\phantomsection\label{sol:m1-2011f-q6-a} 固定前26个人拿到自己的帽子后, 后26人的帽子还可任意排列, 有$26!$种. 因而概率为$26!/52!$; 并没有要求后26人拿错帽子.

\textbf{(b)}\phantomsection\label{sol:m1-2011f-q6-b} 将52人排成序列并按相邻两人配对, 每个无序配对方案因每对内部的两种顺序及26对的排列被数了$2^{26}26!$次. 故配对数为$52!/(2^{26}26!)$. 给定配对后, 每对交换帽子的排列唯一. 除以总排列数得到概率$1/(2^{26}26!)$.
\end{solution}

\originalsection{2012年秋季}
% SOURCE 2012-fall-midterm1
\noindent 官方来源: \href{https://ocw.mit.edu/courses/18-600-probability-and-random-variables-fall-2019/ae6f4359ccee9674fadd34ea6fd23c09_MIT18_600F19_mid1_2012.pdf}{原试卷PDF}; \href{https://ocw.mit.edu/courses/18-600-probability-and-random-variables-fall-2019/d13729835fe69d90045baff75c4ba80c_MIT18_600F19_mid1_2012_soln.pdf}{官方解答PDF}.
\begin{exercise}\label{ex:m1-2012f-q1}
\textbf{2012年秋季, 原题1 (20分).} 独立投掷6枚公平硬币, 设$X$为正面数. 计算:
\begin{enumerate}\item[(a)]\phantomsection\label{ex:m1-2012f-q1-a} $\E[X]$.\item[(b)]\phantomsection\label{ex:m1-2012f-q1-b} $\Var[X]$.\item[(c)]\phantomsection\label{ex:m1-2012f-q1-c} $\Prob(X=6\mid X\ge5)$.\end{enumerate}
\end{exercise}
\begin{solution}
\textbf{(a)}\phantomsection\label{sol:m1-2012f-q1-a} $X$为6个独立$\operatorname{Bernoulli}(1/2)$指标之和, 因而$\E[X]=6/2=3$.

\textbf{(b)}\phantomsection\label{sol:m1-2012f-q1-b} 每个指标方差为$(1/2)(1/2)=1/4$, 独立和的方差为$6/4=3/2$.

\textbf{(c)}\phantomsection\label{sol:m1-2012f-q1-c} $X=6$仅有全正面一种序列, $X=5$有$\binom65=6$种. 条件事件$X\ge5$包含7种等可能序列, 其中仅1种满足$X=6$. 故条件概率为$1/7$.
\end{solution}

\begin{exercise}\label{ex:m1-2012f-q2}
\textbf{2012年秋季, 原题2 (20分).} 4个人各有2顶帽子. 把8顶帽子脱下随机混合, 再分给4个人, 每人仍得到2顶. 假设将8顶不同帽子按每人2顶分配的全部方式等可能.
\begin{enumerate}
\item[(a)]\phantomsection\label{ex:m1-2012f-q2-a} 一共有多少种结果, 即将8顶不同帽子分给4个不同的人且每人2顶的方式数是多少?
\item[(b)]\phantomsection\label{ex:m1-2012f-q2-b} 人标为$1,2,3,4$, 设$E_i$为第$i$个人拿到自己两顶帽子的事件. 计算$\Prob(E_1)$.
\item[(c)]\phantomsection\label{ex:m1-2012f-q2-c} 计算$\Prob(E_1E_2)$, 再计算$\Prob(E_1E_2E_3)$; 后者与$\Prob(E_1E_2E_3E_4)$相同.
\item[(d)]\phantomsection\label{ex:m1-2012f-q2-d} 记$A=\Prob(E_1)$, $B=\Prob(E_1E_2)$, $C=\Prob(E_1E_2E_3)=\Prob(E_1E_2E_3E_4)$. 用容斥把$\Prob(E_1\cup E_2\cup E_3\cup E_4)$表示为$A,B,C$的式子.
\end{enumerate}
\end{exercise}
\begin{solution}
\textbf{(a)}\phantomsection\label{sol:m1-2012f-q2-a} 将8帽排成序列, 相邻2帽依次分给每个人, 每个人内部两帽交换不改变结果. 因而结果数为$8!/(2!)^4=8!/16=2520$.

\textbf{(b)}\phantomsection\label{sol:m1-2012f-q2-b} 第1人的双帽在全部$\binom82$个二元子集中均匀分布, 其中仅1个是他的两顶帽子, 所以$A=1/\binom82=1/28$.

\textbf{(c)}\phantomsection\label{sol:m1-2012f-q2-c} 第1人拿回自己双帽以后, 剩6帽在其余3人间仍按每人2帽均匀分配. 第2人拿回双帽的条件概率为$1/\binom62=1/15$. 因而$B=1/(28\cdot15)=1/420$. 固定前三人的帽子后, 第4人也只能得到自己的帽子, 有利结果仅1种, 故$C=1/2520=16/8!$.

官方解答这一问的中间乘积原印为
\[
\frac18\frac17\,2\,\frac15\frac14\,2
=\frac4{8\cdot7\cdot6\cdot5}=\frac1{420}.
\]
左端第二人的两个因子误印, 应为$1/6$与$1/5$; 原印左端实际等于$1/280$, 不能等于右端. 上述条件计数重新得到的最终$1/420$与官方最终值一致.

\textbf{(d)}\phantomsection\label{sol:m1-2012f-q2-d} 由人之间的对称性, 单事件交有4项, 双事件交有6项, 三事件交有4项, 四事件交有1项. 容斥给出$4A-6B+4C-C=4A-6B+3C$. 若代入精确数值, 为$4/28-6/420+3/2520=109/840$.
\end{solution}

\begin{exercise}\label{ex:m1-2012f-q3}
\textbf{2012年秋季, 原题3 (20分).} 90分钟足球赛中, 每分钟以$2/90$的概率恰进1球, 以$88/90$的概率没有进球, 各分钟相互独立. 设$N$为整场总进球数.
\begin{enumerate}\item[(a)]\phantomsection\label{ex:m1-2012f-q3-a} 精确计算$\E[N]$及$\Var[N]$.\item[(b)]\phantomsection\label{ex:m1-2012f-q3-b} 精确计算恰进1球的概率.\item[(c)]\phantomsection\label{ex:m1-2012f-q3-c} 设$E$为前半场恰0球、后半场恰2球的事件. 用泊松随机变量近似计算$\Prob(E)$.\end{enumerate}
\end{exercise}
\begin{solution}
\textbf{(a)}\phantomsection\label{sol:m1-2012f-q3-a} $N\sim\Bin(90,2/90)$, 因而$\E[N]=90(2/90)=2$, $\Var(N)=90(2/90)(88/90)=176/90=88/45$.

\textbf{(b)}\phantomsection\label{sol:m1-2012f-q3-b} 选出唯一进球的分钟, 所求为$\binom{90}{1}(2/90)(88/90)^{89}$.

\textbf{(c)}\phantomsection\label{sol:m1-2012f-q3-c} 两半场各45分钟, 进球数独立, 各有$\Bin(45,2/90)$分布, 均值为1. 分别用独立$\Pois(1)$随机变量近似, 得
\[
\Prob(E)\approx\left(\ee^{-1}\frac{1^0}{0!}\right)
\left(\ee^{-1}\frac{1^2}{2!}\right)=\frac{\ee^{-2}}2.
\]
官方解答原印第一因子为$\ee^{-\lambda}\lambda^1/1!$, 实际对应1球而非题给0球; 此处应改为$\lambda^0/0!$. 在$\lambda=1$时这两个数偶然相等, 但事件解释仍需修正. 官方PDF还把最终值印为$1/(2\ee^{-2})$, 这大于1, 正确值是$\ee^{-2}/2=1/(2\ee^2)$. 两处均在此明示, 没有把原印值作为有效概率.
\end{solution}

\begin{exercise}\label{ex:m1-2012f-q4}
\textbf{2012年秋季, 原题4 (10分).} 独立掷6个公平六面骰, 每个$6^6$种结果都等可能. 求恰有3个骰子朝上的点数为偶数的概率.
\end{exercise}
\begin{solution}
\phantomsection\label{sol:m1-2012f-q4}
每个骰子出现偶数的概率为$3/6=1/2$, 偶数指标独立. 在6个位置中选3个为偶数, 其余为奇数, 所以概率为$\binom63(1/2)^3(1/2)^3=20/64=5/16$. 骰子有6个点数并不妨碍偶数计数服从参数为$1/2$的二项分布.
\end{solution}

\begin{exercise}\label{ex:m1-2012f-q5}
\textbf{2012年秋季, 原题5 (20分).} 独立掷两个公平六面骰, 第一骰点数为$X$, 第二骰为$Y$, 令$Z=Y+X$. 计算:
\begin{enumerate}\item[(a)]\phantomsection\label{ex:m1-2012f-q5-a} $\E[Z]$.\item[(b)]\phantomsection\label{ex:m1-2012f-q5-b} $\E[Y^2]$.\item[(c)]\phantomsection\label{ex:m1-2012f-q5-c} $\E[YZ]$.\end{enumerate}
\end{exercise}
\begin{solution}
\textbf{(a)}\phantomsection\label{sol:m1-2012f-q5-a} 一个骰子的均值为$(1+2+3+4+5+6)/6=7/2$. 期望线性给出$\E[Z]=\E[X]+\E[Y]=7$.

\textbf{(b)}\phantomsection\label{sol:m1-2012f-q5-b} 六个点数等可能, 故$\E[Y^2]=(1+4+9+16+25+36)/6=91/6$.

\textbf{(c)}\phantomsection\label{sol:m1-2012f-q5-c} 先展开$YZ=Y(Y+X)=Y^2+XY$, 再利用$X,Y$独立, 得$\E[YZ]=\E[Y^2]+\E[X]\E[Y]=91/6+49/4=329/12$. 官方解答左端曾漏掉期望符号, 原印为$[YZ]$; 上式明确其含义为$\E[YZ]$.
\end{solution}

\begin{exercise}\label{ex:m1-2012f-q6}
\textbf{2012年秋季, 原题6 (10分).} 无限次独立掷一枚正面概率为$1/3$的硬币, 第3次正面出现于第$X$次投掷.
\begin{enumerate}\item[(a)]\phantomsection\label{ex:m1-2012f-q6-a} 计算$\Prob(X=9)$.\item[(b)]\phantomsection\label{ex:m1-2012f-q6-b} 计算$\E[X]$.\end{enumerate}
\end{exercise}
\begin{solution}
\textbf{(a)}\phantomsection\label{sol:m1-2012f-q6-a} 第9次须为正面, 前8次恰有2个正面, 因而$\Prob(X=9)=\binom82(1/3)^2(2/3)^6(1/3)=\binom82(1/3)^3(2/3)^6$.

\textbf{(b)}\phantomsection\label{sol:m1-2012f-q6-b} $X$由三个独立的等待正面段长相加, 每段为$\Geom(1/3)$, 均值为3. 期望可加得$\E[X]=3\cdot3=9$. 分段独立性的依据与习题\ref{ex:m1-2011f-q4}相同.
\end{solution}

\originalsection{2014年春季}
% SOURCE 2014-spring-midterm1
\noindent 官方来源: \href{https://ocw.mit.edu/courses/18-600-probability-and-random-variables-fall-2019/d88a1f88b725d3ec75144e5dd130fd55_MIT18_600F19_mid1_2014.pdf}{原试卷PDF}; \href{https://ocw.mit.edu/courses/18-600-probability-and-random-variables-fall-2019/1c1579e396bb44c47ca34dd6f5f2e6ba_MIT18_600F19_mid1_2014_soln.pdf}{官方解答PDF}.
\begin{exercise}\label{ex:m1-2014s-q1}
\textbf{2014年春季, 原题1 (10分).} 有多少个非负整数五元组$(a_1,a_2,a_3,a_4,a_5)$满足$a_1+a_2+a_3+a_4+a_5=100$?
\end{exercise}
\begin{solution}
\phantomsection\label{sol:m1-2014s-q1}
用100颗星和4条分隔线把五个分量依次编码, 允许空段以表示零. 每个解与104个位置中选择4个分隔线位置一一对应, 所以答案为$\binom{104}{4}=104!/(4!100!)$.
\end{solution}

\begin{exercise}\label{ex:m1-2014s-q2}
\textbf{2014年春季, 原题2 (20分).} 30个人受邀参加聚会. 每人独立地以$1/3$的概率接受邀请. 设$X$为接受邀请的人数. 计算:
\begin{enumerate}\item[(a)]\phantomsection\label{ex:m1-2014s-q2-a} $\E[X]$.\item[(b)]\phantomsection\label{ex:m1-2014s-q2-b} $\Var[X]$.\item[(c)]\phantomsection\label{ex:m1-2014s-q2-c} $\E[X^2]$.\item[(d)]\phantomsection\label{ex:m1-2014s-q2-d} $\E[X^2-4X+5]$.\end{enumerate}
\end{exercise}
\begin{solution}
\textbf{(a)}\phantomsection\label{sol:m1-2014s-q2-a} $X\sim\Bin(30,1/3)$, 可写为30个独立接受邀请指标之和, 故$\E[X]=30/3=10$.

\textbf{(b)}\phantomsection\label{sol:m1-2014s-q2-b} 每个指标方差为$(1/3)(2/3)$, 所以$\Var(X)=30(1/3)(2/3)=20/3$.

\textbf{(c)}\phantomsection\label{sol:m1-2014s-q2-c} 从方差恒等式恢复二阶矩: $\E[X^2]=\Var(X)+\E[X]^2=20/3+100=320/3$.

\textbf{(d)}\phantomsection\label{sol:m1-2014s-q2-d} 按期望线性逐项取期望, 得$\E[X^2-4X+5]=320/3-4\cdot10+5=215/3$. 这里不用求$X^2-4X+5$的完整分布.
\end{solution}

\begin{exercise}\label{ex:m1-2014s-q3}
\textbf{2014年春季, 原题3 (20分).} Bob注意到他的干洗店Facebook页面, 每分钟独立地以$1/720$的概率获得一个赞. 设$L$为24小时内收到的总赞数.
\begin{enumerate}\item[(a)]\phantomsection\label{ex:m1-2014s-q3-a} 精确计算$\E[L]$及$\Var[L]$.\item[(b)]\phantomsection\label{ex:m1-2014s-q3-b} 精确计算$\Prob(L=0)$.\item[(c)]\phantomsection\label{ex:m1-2014s-q3-c} Bob希望在接下来的24小时内至少多得到2个赞, 这会让累计赞数达到三位数. 用泊松随机变量近似计算$\Prob(L\ge2)$.\end{enumerate}
\end{exercise}
\begin{solution}
\textbf{(a)}\phantomsection\label{sol:m1-2014s-q3-a} 24小时共有$24\cdot60=1440$分钟, 因而$L\sim\Bin(1440,1/720)$. 独立指标求和给出$\E[L]=2$, $\Var(L)=2(719/720)=719/360$. 官方PDF所写为$2\frac{719}{720}$, 文本提取出现的“2719”是乘数与分子粘连, 不应解读成分子2719.

\textbf{(b)}\phantomsection\label{sol:m1-2014s-q3-b} 所有1440分钟均未获赞的概率为$(719/720)^{1440}$.

\textbf{(c)}\phantomsection\label{sol:m1-2014s-q3-c} 用均值$\lambda=2$的泊松分布近似, 再减去0个和1个赞的概率, 得$\Prob(L\ge2)\approx1-\ee^{-2}-2\ee^{-2}=1-3\ee^{-2}$. 官方解答在此原印“approximately binomial with parameter $\lambda$”, 分布名称应为Poisson; $L$本来精确地服从二项分布, 此处近似对象才是泊松分布. 最终近似值无需改变.
\end{solution}

\begin{exercise}\label{ex:m1-2014s-q4}
\textbf{2014年春季, 原题4 (10分).} 无限次独立掷一枚正面概率为$p$的硬币. 用$p$表示第5次正面恰出现在第10次投掷的概率.
\end{exercise}
\begin{solution}
\phantomsection\label{sol:m1-2014s-q4}
本题与习题\ref{ex:m1-2011s-q1}的(c)条件及所问完全相同. 其事件是前9次恰4次正面且第10次为正面, 两部分独立, 故仍为$\binom94p^5(1-p)^5$. 前题已说明位置计数及端点情形.
\end{solution}

\begin{exercise}\label{ex:m1-2014s-q5}
\textbf{2014年春季, 原题5 (20分).} $X$为标准公平六面骰的点数, 在$\{1,2,3,4,5,6\}$中等概率取值. $Y$为同一骰子另一次独立投掷的点数. 计算:
\begin{enumerate}\item[(a)]\phantomsection\label{ex:m1-2014s-q5-a} $\E[X^2]$.\item[(b)]\phantomsection\label{ex:m1-2014s-q5-b} $\E[XY]$.\item[(c)]\phantomsection\label{ex:m1-2014s-q5-c} $\Cov(X,Y)=\E[XY]-\E[X]\E[Y]$.\end{enumerate}
\end{exercise}
\begin{solution}
\textbf{(a)}\phantomsection\label{sol:m1-2014s-q5-a} 与习题\ref{ex:m1-2012f-q5}的(b)仅骰子命名不同, 同样的六点均匀平均给出$\E[X^2]=91/6$.

\textbf{(b)}\phantomsection\label{sol:m1-2014s-q5-b} 独立性使联合概率分解, 故
\[
\E[XY]=\sum_{x=1}^6\sum_{y=1}^6xy\frac1{36}
=\left(\sum_{x=1}^6\frac x6\right)
\left(\sum_{y=1}^6\frac y6\right)
=\frac{49}{4}.
\]
\textbf{(c)}\phantomsection\label{sol:m1-2014s-q5-c} $\E[X]=\E[Y]=7/2$, 把(b)代入协方差定义得$\Cov(X,Y)=49/4-(7/2)^2=0$. 本题的独立性是零协方差的充分条件.
\end{solution}

\begin{exercise}\label{ex:m1-2014s-q6}
\textbf{2014年春季, 原题6 (20分).} 3顶帽子从各自指定的箱中掉出, 再随机放回, 每箱一顶, 全部$3!$种分配等可能. 计算:
\begin{enumerate}
\item[(a)]\phantomsection\label{ex:m1-2014s-q6-a} 回到自己箱中的帽子数的期望.
\item[(b)]\phantomsection\label{ex:m1-2014s-q6-b1} 第3顶帽子回到自己箱中的概率.
\item[(b)]\phantomsection\label{ex:m1-2014s-q6-b2} 在第1顶帽子未回到自己箱中的条件下, 第3顶帽子回到自己箱中的条件概率. \textnormal{(原题及官方解答均重复印为(b); 此项为第二个(b).)}
\end{enumerate}
\end{exercise}
\begin{solution}
\textbf{(a)}\phantomsection\label{sol:m1-2014s-q6-a} 令$I_i$表示第$i$顶帽子归位. 每帽在3个箱中等可能, 所以$\E[I_i]=1/3$. 归位数$I_1+I_2+I_3$的期望为$3/3=1$.

\textbf{(b), 第一次}\phantomsection\label{sol:m1-2014s-q6-b1} 第3帽在3箱中的位置均匀, 归位概率为$1/3$.

\textbf{(b), 第二次}\phantomsection\label{sol:m1-2014s-q6-b2} 令$A$为第3帽归位, $B$为第1帽未归位, 则$\Prob(B)=2/3$. 第3帽归位后, 前两帽只有原位和交换两种分配, 仅交换满足$B$. 因而在全部6种排列中$AB$只有1种, $\Prob(AB)=1/6$. 条件概率为$\Prob(A\mid B)=(1/6)/(2/3)=1/4$. 第1帽未归位使归位事件之间的关系改变, 不能把无条件$1/3$直接沿用.
\end{solution}

\originalsection{2016年春季}
% SOURCE 2016-spring-midterm1
\noindent 官方来源: \href{https://ocw.mit.edu/courses/18-600-probability-and-random-variables-fall-2019/b146bf10509052fa0881160aaa603aea_MIT18_600F19_mid1_2016.pdf}{原试卷PDF}; \href{https://ocw.mit.edu/courses/18-600-probability-and-random-variables-fall-2019/b744aee08a1ab02c622fc9916e4c7467_MIT18_600F19_mid1_2016_soln.pdf}{官方解答PDF}.
\begin{exercise}\label{ex:m1-2016s-q1}
\textbf{2016年春季, 原题1 (20分).} $X_1,X_2,X_3$相互独立, 各以$1/4$概率取1、以$3/4$概率取0. 令$Y=X_1+X_2$, $Z=X_2+X_3$. 计算:
\begin{enumerate}\item[(a)]\phantomsection\label{ex:m1-2016s-q1-a} $\E[Y]$.\item[(b)]\phantomsection\label{ex:m1-2016s-q1-b} $\E[YZ]$.\item[(c)]\phantomsection\label{ex:m1-2016s-q1-c} $\E[Y^2]$.\item[(d)]\phantomsection\label{ex:m1-2016s-q1-d} $\Var[Y]$.\item[(e)]\phantomsection\label{ex:m1-2016s-q1-e} $\Cov(Y,Z)$. 提醒: $\Cov(Y,Z)=\E[YZ]-\E[Y]\E[Z]$.\end{enumerate}
\end{exercise}
\begin{solution}
每个$X_i$均值为$1/4$, 且$X_i^2=X_i$. 不同指标的乘积期望因独立性而为$1/16$; 同一个指标的平方期望为$1/4$, 这一区分是计算混合矩的要点.

\textbf{(a)}\phantomsection\label{sol:m1-2016s-q1-a} 期望可加给出$\E[Y]=1/4+1/4=1/2$, 同样$\E[Z]=1/2$.

\textbf{(b)}\phantomsection\label{sol:m1-2016s-q1-b} 展开$YZ=X_1X_2+X_1X_3+X_2^2+X_2X_3$, 所以$\E[YZ]=1/16+1/16+1/4+1/16=7/16$.

\textbf{(c)}\phantomsection\label{sol:m1-2016s-q1-c} 展开$Y^2=X_1^2+2X_1X_2+X_2^2$, 得$\E[Y^2]=1/4+2/16+1/4=5/8$.

\textbf{(d)}\phantomsection\label{sol:m1-2016s-q1-d} $\Var(Y)=5/8-(1/2)^2=3/8$. 也可从$Y\sim\Bin(2,1/4)$求得$2(1/4)(3/4)=3/8$.

\textbf{(e)}\phantomsection\label{sol:m1-2016s-q1-e} $\Cov(Y,Z)=7/16-1/4=3/16$. 两个和共享$X_2$, 并不独立. 另一种看法是协方差对两个变量分别线性, 只有共享项留下$\Cov(X_2,X_2)=\Var(X_2)=3/16$, 其余独立项协方差为零.
\end{solution}

\begin{exercise}\label{ex:m1-2016s-q2}
\textbf{2016年春季, 原题2 (20分).} 某小型学院的新生班有400名学生. 每个学生独立地以$1/200$的概率参加击剑队选拔. 设$N$为参加选拔的新生总数.
\begin{enumerate}\item[(a)]\phantomsection\label{ex:m1-2016s-q2-a} 精确计算$\E[N]$.\item[(b)]\phantomsection\label{ex:m1-2016s-q2-b} 精确计算$\Var[N]$.\item[(c)]\phantomsection\label{ex:m1-2016s-q2-c} 对$k\in\{0,1,\ldots,400\}$写出$\Prob(N=k)$的精确公式.\item[(d)]\phantomsection\label{ex:m1-2016s-q2-d} 用泊松随机变量近似计算$\Prob(N=3)$.\end{enumerate}
\end{exercise}
\begin{solution}
\textbf{(a)}\phantomsection\label{sol:m1-2016s-q2-a} $N$是400个独立参加指标之和, 所以$\E[N]=400/200=2$.

\textbf{(b)}\phantomsection\label{sol:m1-2016s-q2-b} 方差按独立性可加, $\Var(N)=400(1/200)(199/200)=199/100=1.99$, 有限小数在此是精确值.

\textbf{(c)}\phantomsection\label{sol:m1-2016s-q2-c} 从400人中选$k$人参加, 各指定参加集合的概率为$(1/200)^k(199/200)^{400-k}$. 因而$\Prob(N=k)=\binom{400}{k}(1/200)^k(199/200)^{400-k}$, 该公式包括$k=0,400$.

\textbf{(d)}\phantomsection\label{sol:m1-2016s-q2-d} 用均值2的$\Pois(2)$近似, 得$\Prob(N=3)\approx\ee^{-2}2^3/3!=4\ee^{-2}/3$. 同均值并不意味着方差精确相同, 此处只按题目要求做小概率二项的泊松近似.
\end{solution}

\begin{exercise}\label{ex:m1-2016s-q3}
\textbf{2016年春季, 原题3 (15分).} 一副牌有52种不同牌, 每种恰1张. 现有两副相同的牌, 共104张, 意外丢失10张, 丢失集合在全部10张子集中均匀选择, 剩下94张. 求能从剩牌组成完整一副牌的概率, 即剩牌仍包含52种中的每一种. 提示: 等价于丢失的10张中每种牌至多出现1张.
\end{exercise}
\begin{solution}
\phantomsection\label{sol:m1-2016s-q3}
虽然同种牌来自两副而外观相同, 它们是抽样模型中的两张不同实体; 全部丢失集合有$\binom{104}{10}$种. 剩牌可组成完整牌组当且仅当没有任何牌种两张全丢, 即10张失牌来自10个不同牌种. 先选10种牌, 有$\binom{52}{10}$种, 再对每种选来自哪一副的那张, 有$2^{10}$种. 这些选择唯一确定有利失牌集合, 所以概率为$2^{10}\binom{52}{10}/\binom{104}{10}$.
\end{solution}

\begin{exercise}\label{ex:m1-2016s-q4}
\textbf{2016年春季, 原题4 (15分).} 7个人将帽子投入箱中随机混合再每人拿一顶. 设$N$为拿回自己帽子的人数. 计算:
\begin{enumerate}\item[(a)]\phantomsection\label{ex:m1-2016s-q4-a} $\E[N]$.\item[(b)]\phantomsection\label{ex:m1-2016s-q4-b} $\Prob(N=5)$.\item[(c)]\phantomsection\label{ex:m1-2016s-q4-c} $\Prob(N=7\mid N\ge5)$.\end{enumerate}
\end{exercise}
\begin{solution}
\textbf{(a)}\phantomsection\label{sol:m1-2016s-q4-a} 每人拿回自己帽子的概率为$1/7$, 加总7个指标的均值得$\E[N]=1$.

\textbf{(b)}\phantomsection\label{sol:m1-2016s-q4-b} 恰5人归位时, 剩下2人的帽子必须交换. 选择这2人有$\binom72$种, 给定2人后交换唯一. 故$\Prob(N=5)=\binom72/7!$.

\textbf{(c)}\phantomsection\label{sol:m1-2016s-q4-c} 不可能恰6人归位: 6顶已被其主人拿走后, 第7人只能拿剩余自己的帽子. 因而$N\ge5$只有$N=5$和$N=7$两种情况, $N=7$对应唯一恒等排列. 条件概率为
\[
\frac{1/7!}{\binom72/7!+1/7!}=\frac1{21+1}=\frac1{22}.
\]
分母若误加$N=6$的虚构排列就会破坏此条件计数.
\end{solution}

\begin{exercise}\label{ex:m1-2016s-q5}
\textbf{2016年春季, 原题5 (10分).}罐中有10个黑球和10个白球. 从中均匀随机选出8球, 求恰有4黑4白的概率.
\end{exercise}
\begin{solution}
\phantomsection\label{sol:m1-2016s-q5}
全部8球子集有$\binom{20}{8}$种. 有利子集由黑球中选4球及白球中选4球唯一确定, 有$\binom{10}{4}\binom{10}{4}$种. 无放回取样下所求概率为$\binom{10}{4}^2/\binom{20}{8}$, 不能使用独立取样的二项式概率.
\end{solution}

\begin{exercise}\label{ex:m1-2016s-q6}
\textbf{2016年春季, 原题6 (20分).} $X$是标准公平六面骰的点数, 在$\{1,2,3,4,5,6\}$中等概率取值; $Y$为同一骰子的独立另一次投掷的点数. 计算:
\begin{enumerate}\item[(a)]\phantomsection\label{ex:m1-2016s-q6-a} $\E[X]$.\item[(b)]\phantomsection\label{ex:m1-2016s-q6-b} $\E[X^2]$.\item[(c)]\phantomsection\label{ex:m1-2016s-q6-c} $\E[5X^7-5Y^7+5]$.\item[(d)]\phantomsection\label{ex:m1-2016s-q6-d} $\E[XY]$.\end{enumerate}
\end{exercise}
\begin{solution}
\textbf{(a)}\phantomsection\label{sol:m1-2016s-q6-a} 均匀取值给出$\E[X]=(1+2+3+4+5+6)/6=7/2$.

\textbf{(b)}\phantomsection\label{sol:m1-2016s-q6-b} 与习题\ref{ex:m1-2012f-q5}的(b)及习题\ref{ex:m1-2014s-q5}的(a)相同, 二阶矩为$91/6$.

\textbf{(c)}\phantomsection\label{sol:m1-2016s-q6-c} 因为$X,Y$同分布, $\E[X^7]=\E[Y^7]$, 两者均有限. 期望线性使高次项抵消, 得$5\E[X^7]-5\E[Y^7]+5=5$. 这一抵消只需同分布, 此步骤本身不依赖独立性.

\textbf{(d)}\phantomsection\label{sol:m1-2016s-q6-d} 与习题\ref{ex:m1-2014s-q5}的(b)完全相同, 独立性给出$\E[XY]=\E[X]\E[Y]=(7/2)^2=49/4$; 该处已展开联合概率的乘积推导.
\end{solution}

\originalsection{2017年春季}
% SOURCE 2017-spring-midterm1
\noindent 官方来源: \href{https://ocw.mit.edu/courses/18-600-probability-and-random-variables-fall-2019/25bb4897aed3cc4112d685e4dfbb9a7d_MIT18_600F19_mid1_2017.pdf}{原试卷PDF}; \href{https://ocw.mit.edu/courses/18-600-probability-and-random-variables-fall-2019/87a8eb5f64734f2b436ac96d15cfe6bf_MIT18_600F19_mid1_2017_soln.pdf}{官方解答PDF}.
\begin{exercise}\label{ex:m1-2017s-q1}
\textbf{2017年春季, 原题1 (20分).} $X$为标准公平六面骰的点数, 在$\{1,2,3,4,5,6\}$中等概率取值. $Y,Z$为另两次独立投掷的点数, 即$X,Y,Z$相互独立. 计算:
\begin{enumerate}\item[(a)]\phantomsection\label{ex:m1-2017s-q1-a} $\E[X]$及$\E[X^2]$.\item[(b)]\phantomsection\label{ex:m1-2017s-q1-b} $\Var[X]$.\item[(c)]\phantomsection\label{ex:m1-2017s-q1-c} $\E[X^2+2Y^2+3Z^2]$.\item[(d)]\phantomsection\label{ex:m1-2017s-q1-d} 条件概率$\Prob(X=1\mid X=Y)$.\end{enumerate}
\end{exercise}
\begin{solution}
\textbf{(a)}\phantomsection\label{sol:m1-2017s-q1-a} 对六个等可能点数直接平均, $\E[X]=21/6=7/2$, $\E[X^2]=91/6$. 本小问同时要求两个矩, 两项都保留.

\textbf{(b)}\phantomsection\label{sol:m1-2017s-q1-b} $\Var(X)=\E[X^2]-\E[X]^2=91/6-49/4=(182-147)/12=35/12$.

\textbf{(c)}\phantomsection\label{sol:m1-2017s-q1-c} 三骰同分布, 用期望线性得$\E[X^2]+2\E[Y^2]+3\E[Z^2]=(1+2+3)91/6=91$. 官方解答展开式最后一项原印$3[Z^2]$, 漏了期望符号; 正确写法为$3\E[Z^2]$.

\textbf{(d)}\phantomsection\label{sol:m1-2017s-q1-d} $X=Y$对应$36$种等可能双骰结果中的$6$种, 概率$1/6$; $X=1$且$X=Y$仅$(1,1)$一种, 概率$1/36$. 因而条件概率为$(1/36)/(1/6)=1/6$. 第三骰的值没有进入这两个事件, 其概率质量在边缘化时加总为1.
\end{solution}

\begin{exercise}\label{ex:m1-2017s-q2}
\textbf{2017年春季, 原题2 (10分).} 罐中有黑、白、红、黄球各5个. 均匀随机选择8球, 求每种颜色恰有2球的概率.
\end{exercise}
\begin{solution}
\phantomsection\label{sol:m1-2017s-q2}
20个球的8球子集共有$\binom{20}{8}$种. 从每一颜色的5个球中选2球, 四个选择独立地确定一个有利子集, 因而有利子集数为$\binom52^4$. 概率为$\binom52^4/\binom{20}{8}$. 这里“选择独立地确定”是组合乘法原理, 并非说无放回后的颜色计数随机变量相互独立.
\end{solution}

\begin{exercise}\label{ex:m1-2017s-q3}
\textbf{2017年春季, 原题3 (20分).} $M$表示一个人有某种疾病, $T$表示疾病检测呈阳性. 此人患病概率为$1/20$; 已患病时检测以90\%概率正确识别疾病; 未患病时也有20\%概率呈阳性. 即$\Prob(M)=1/20$, $\Prob(T\mid M)=9/10$, $\Prob(T\mid M^c)=1/5$. 计算:
\begin{enumerate}\item[(a)]\phantomsection\label{ex:m1-2017s-q3-a} $\Prob(T)$, 即检测总体阳性概率.\item[(b)]\phantomsection\label{ex:m1-2017s-q3-b} $\Prob(M\mid T)$, 即检测阳性后患病概率.\item[(c)]\phantomsection\label{ex:m1-2017s-q3-c} $\Prob(M\mid T^c)$, 即检测阴性后患病概率.\end{enumerate}
\end{exercise}
\begin{solution}
\textbf{(a)}\phantomsection\label{sol:m1-2017s-q3-a} 按患病与未患病这两个互斥完备事件分解阳性概率,
\[
\Prob(T)=\frac1{20}\frac9{10}+\frac{19}{20}\frac15
=\frac9{200}+\frac{38}{200}=\frac{47}{200}.
\]
\textbf{(b)}\phantomsection\label{sol:m1-2017s-q3-b} 患病且阳性的概率为$9/200$, 除以上述阳性总概率, 得$\Prob(M\mid T)=(9/200)/(47/200)=9/47$. 即使检测对患者有90\%的识别率, 阳性人群中也包含人数更多的未患者, 不能把条件方向反过来.

\textbf{(c)}\phantomsection\label{sol:m1-2017s-q3-c} 患病且阴性的概率为$(1/20)(1/10)=1/200$, 阴性总概率为$1-47/200=153/200$. 因而$\Prob(M\mid T^c)=1/153$. 两个分母均为正, 条件概率有定义; 本题只计算给定的历史教学模型.
\end{solution}

\begin{exercise}\label{ex:m1-2017s-q4}
\textbf{2017年春季, 原题4 (10分).} 标准52张牌分4种花色, 每色13张. 选取无序13张牌集合共有$\binom{52}{13}$种. 有多少个无序13张牌集合, 其中恰7张属于同一花色, 另6张不属于该花色?
\end{exercise}
\begin{solution}
\phantomsection\label{sol:m1-2017s-q4}
先选该花色, 有4种选择; 再从该花色13张中选7张, 有$\binom{13}{7}$种; 从其他39张选6张, 有$\binom{39}{6}$种. 这里不会重复计数: 若两种不同花色都出现7张, 至少需要14张牌, 超过题给13张. 所以指定的7张花色在每个有利集合中唯一, 答案为$4\binom{13}{7}\binom{39}{6}$.
\end{solution}

\begin{exercise}\label{ex:m1-2017s-q5}
\textbf{2017年春季, 原题5 (10分).} 聚会有12片披萨和4个人, 将这些片分给4个人有多少种方式? 即有多少个非负整数有序四元组$(a_1,a_2,a_3,a_4)$满足$a_1+a_2+a_3+a_4=12$?
\end{exercise}
\begin{solution}
\phantomsection\label{sol:m1-2017s-q5}
题目的整数四元组表述按每人获得的片数区分结果, 不按每片实体区分. 用12颗星和3条分隔线编码四个分量, 共15个位置, 所以解数为$\binom{15}{3}=455$. 允许某人没有披萨对应于空段.
\end{solution}

\begin{exercise}\label{ex:m1-2017s-q6}
\textbf{2017年春季, 原题6 (20分).} Harry是一位水平较弱的篮球运动员, 正在练习罚球. 每次出手以$1/100$的概率命中, 各次相互独立. 设$X$为100次出手后的命中数.
\begin{enumerate}
\item[(a)]\phantomsection\label{ex:m1-2017s-q6-a} 精确计算$\Var[X]$.
\item[(b)]\phantomsection\label{ex:m1-2017s-q6-b} 用泊松近似估计$\Prob(X=2)$.
\item[(c)]\phantomsection\label{ex:m1-2017s-q6-c} Harry希望$X>1$发生, 这样命中数会因好运而超过他的期望值, 或许会给在场球探留下印象. 用泊松近似估计$\Prob(X>1)$.
\item[(d)]\phantomsection\label{ex:m1-2017s-q6-d} 用同一泊松近似估计$\Prob(X<1)$. 它比(c)中的估计值大、小还是相同?
\end{enumerate}
\end{exercise}
\begin{solution}
\textbf{(a)}\phantomsection\label{sol:m1-2017s-q6-a} $X\sim\Bin(100,1/100)$, $\E[X]=1$, 独立指标方差相加得$\Var(X)=100(1/100)(99/100)=99/100$.

\textbf{(b)}\phantomsection\label{sol:m1-2017s-q6-b} 用$\Pois(1)$近似, 得$\Prob(X=2)\approx\ee^{-1}1^2/2!=1/(2\ee)$.

\textbf{(c)}\phantomsection\label{sol:m1-2017s-q6-c} 由于$X$为整数, $X>1$即至少2次命中. 在泊松近似下, $\Prob(X=0)\approx1/\ee$, $\Prob(X=1)\approx1/\ee$, 所以$\Prob(X>1)\approx1-2/\ee$.

\textbf{(d)}\phantomsection\label{sol:m1-2017s-q6-d} 非负整数$X<1$等价于$X=0$, 故近似值为$1/\ee$. 比较差值$1/\ee-(1-2/\ee)=3/\ee-1$. 由
\[
\ee=\sum_{k=0}^{\infty}\frac1{k!}<1+1+\sum_{k=2}^{\infty}\frac1{2^{k-1}}=3
\]
可知差值为正, 因而$\Prob(X<1)$的泊松估计大于$\Prob(X>1)$的估计. 严格不等号来自$k\ge3$时$k!>2^{k-1}$. 两侧事件虽围绕均值1, 分布并不对称; 少于均值只能差1次, 超过均值则可多出许多次.
\end{solution}

\begin{exercise}\label{ex:m1-2017s-q7}
\textbf{2017年春季, 原题7 (10分).} 固定整数$n\ge3$. 从$\{1,2,\ldots,n\}$的$n!$种排列中均匀选择$\sigma$. 设$X=X(\sigma)$为$\sigma$中长度3的循环数, 计算$\E[X]$.

如有帮助可用以下记号: 对$i<j<k$, $E_{i,j,k}$为这三个数恰组成一个长度3循环的事件, 即或者$\sigma(i)=j,\sigma(j)=k,\sigma(k)=i$, 或者$\sigma(i)=k,\sigma(k)=j,\sigma(j)=i$. 令$X_{i,j,k}=\ind_{E_{i,j,k}}$.
\end{exercise}
\begin{solution}
\phantomsection\label{sol:m1-2017s-q7}
对给定无序三元组$\{i,j,k\}$, 恰有两种循环方向. 固定一种方向后, 其余$n-3$个元素可任意排列, 有$(n-3)!$种. 因而
\[
\E[X_{i,j,k}]=\Prob(E_{i,j,k})=\frac{2(n-3)!}{n!}
=\frac2{n(n-1)(n-2)}.
\]
每个长度3循环有唯一一个无序三元组作为支撑, 所以$X=\sum_{i<j<k}X_{i,j,k}$; 这里不按循环的起点重复计数. 有$\binom n3$个三元组, 期望线性给出$\E[X]=\binom n3\,2/[n(n-1)(n-2)]=1/3$. 指标之间是否独立不影响期望. 当$n=3$时, 两个三循环占6个排列中的2个, 也得到$1/3$.
\end{solution}

\originalsection{2018年春季}
% SOURCE 2018-spring-midterm1
\noindent 官方来源: \href{https://ocw.mit.edu/courses/18-600-probability-and-random-variables-fall-2019/0a4597f4544a9758f68ad1f5441f8233_MIT18_600F19_mid1_2018.pdf}{原试卷PDF}; \href{https://ocw.mit.edu/courses/18-600-probability-and-random-variables-fall-2019/134c16cb0b85b1dd5ff8fad6383b3415_MIT18_600F19_mid1_2018_soln.pdf}{官方解答PDF}.
\begin{exercise}\label{ex:m1-2018s-q1}
\textbf{2018年春季, 原题1 (20分).} 独立掷6个标准公平六面骰, 设$X$为点数显示6的骰子数.
\begin{enumerate}\item[(a)]\phantomsection\label{ex:m1-2018s-q1-a} 计算$\E[X]$.\item[(b)]\phantomsection\label{ex:m1-2018s-q1-b} 计算$\Var(X)$.\item[(c)]\phantomsection\label{ex:m1-2018s-q1-c} 计算$\Prob(X=5)$.\item[(d)]\phantomsection\label{ex:m1-2018s-q1-d} 计算$\Prob(X=6\mid X\ge5)$.\end{enumerate}
\end{exercise}
\begin{solution}
\textbf{(a)}\phantomsection\label{sol:m1-2018s-q1-a} 每个显示6的指标成功率为$1/6$, 且指标独立, 故$X\sim\Bin(6,1/6)$, $\E[X]=6(1/6)=1$. 官方解答把计数对象误称为“正面数”, 此处依题面明确为显示6的骰子数; 二项参数及答案一致.

\textbf{(b)}\phantomsection\label{sol:m1-2018s-q1-b} 独立指标的方差相加, $\Var(X)=6(1/6)(5/6)=5/6$.

\textbf{(c)}\phantomsection\label{sol:m1-2018s-q1-c} 在6个位置中选5个显示6, 余下一个有5种非6的点数. 所以概率为$\binom65(1/6)^5(5/6)=30/6^6=5/6^5$.

\textbf{(d)}\phantomsection\label{sol:m1-2018s-q1-d} 全部为6仅1种序列, 恰5个为6有30种序列, 均来自$6^6$个等可能结果. 因而
\[
\Prob(X=6\mid X\ge5)=\frac{1/6^6}{(1+30)/6^6}=\frac1{31}.
\]
同为6次试验, 本题成功率$1/6$与2012年公平硬币的$1/2$不同, 条件概率也随之改变.
\end{solution}

\begin{exercise}\label{ex:m1-2018s-q2}
\textbf{2018年春季, 原题2 (10分).} 三个事件$E,F,G$满足$\Prob(E)=\Prob(F)=\Prob(G)=0.4$, $\Prob(EF)=\Prob(EG)=\Prob(FG)=0.2$, 且$\Prob(EFG)=0.1$. 计算$\Prob(E\cup F\cup G)$.
\end{exercise}
\begin{solution}
\phantomsection\label{sol:m1-2018s-q2}
单个概率相加把每个双交中的点多算一次, 因而减去三个双交; 三交中的点先被加3次又减3次, 还需加回1次. 三事件容斥给出$\Prob(E\cup F\cup G)=3(0.4)-3(0.2)+0.1=0.7$. 本题没有给独立性, 也不需要用独立性推测交概率.
\end{solution}

\begin{exercise}\label{ex:m1-2018s-q3}
\textbf{2018年春季, 原题3 (10分).} 计算:
\begin{enumerate}\item[(a)]\phantomsection\label{ex:m1-2018s-q3-a} $\displaystyle\lim_{n\to\infty}(1+5/n)^n$.\item[(b)]\phantomsection\label{ex:m1-2018s-q3-b} $\displaystyle\sum_{n=0}^{\infty}5^n/n!$.\end{enumerate}
\end{exercise}
\begin{solution}
\textbf{(a)}\phantomsection\label{sol:m1-2018s-q3-a} 用基本极限$\ln(1+t)/t\to1$当$t\to0$, 得$n\ln(1+5/n)=5\ln(1+5/n)/(5/n)\to5$. 指数函数连续, 故原式趋于$\ee^5$. 该步骤作用于所有正整数$n$, 无须把$n$限制为5的倍数.

\textbf{(b)}\phantomsection\label{sol:m1-2018s-q3-b} 指数函数的展开为$\ee^x=\sum_{n\ge0}x^n/n!$. 它在$x=5$仍收敛, 因为相邻项之比为$5/(n+1)\to0$. 代入$x=5$得到$\sum_{n\ge0}5^n/n!=\ee^5$. 极限和级数形式不同, 这里都表达指数函数在同一参数的值.
\end{solution}

\begin{exercise}\label{ex:m1-2018s-q4}
\textbf{2018年春季, 原题4 (20分).} Alice的新餐馆营业头几个月有10000人到访. 每人独立地以$5/10000$的概率留下Yelp好评, 如“萨莫萨很好吃!”, 以$1/10000$的概率留下差评, 如“服务员很粗鲁!”, 以$9994/10000$的概率不留评价. 设$X$为好评总数, $Y$为差评总数.
\begin{enumerate}\item[(a)]\phantomsection\label{ex:m1-2018s-q4-a} 精确计算$\E[Y]$及$\Var(Y)$.\item[(b)]\phantomsection\label{ex:m1-2018s-q4-b} 用泊松随机变量近似计算$\Prob(X=3)$.\item[(c)]\phantomsection\label{ex:m1-2018s-q4-c} 用泊松随机变量近似计算$\Prob(Y=0)$.\end{enumerate}
\end{exercise}
\begin{solution}
每位顾客的好评与差评互斥, 因而$X,Y$不要求相互独立; 但不同顾客的行为独立, 所以两个单独计数都各有二项边缘分布.

\textbf{(a)}\phantomsection\label{sol:m1-2018s-q4-a} 差评指标成功率为$1/10000$, $Y\sim\Bin(10000,1/10000)$. 因而$\E[Y]=1$, 而方差为
\[
\Var(Y)=10000\frac1{10000}\frac{9999}{10000}=\frac{9999}{10000}.
\]

\textbf{(b)}\phantomsection\label{sol:m1-2018s-q4-b} $X\sim\Bin(10000,5/10000)$, 均值为5, 用$\Pois(5)$近似得$\Prob(X=3)\approx\ee^{-5}5^3/3!=125/(6\ee^5)$.

\textbf{(c)}\phantomsection\label{sol:m1-2018s-q4-c} 用$\Pois(1)$近似$Y$, 零次的概率为$\ee^{-1}1^0/0!=1/\ee$. 官方解答原印“$k=1$且$\lambda=1$”, 对零次事件应为$k=0$. 当$\lambda=1$时零次和一次的泊松概率恰相同, 故最终$1/\ee$不变, 但所代入的事件参数必须纠正. 精确零次概率为$(1-1/10000)^{10000}$, 与上述近似的区别在此明确.
\end{solution}

\begin{exercise}\label{ex:m1-2018s-q5}
\textbf{2018年春季, 原题5 (10分).} 一副牌有120张: 红牌30张、黑牌40张、蓝牌50张. 均匀随机选12张, 全部12张子集等可能. 求所选集合恰有3红、4黑、5蓝的概率.
\end{exercise}
\begin{solution}
\phantomsection\label{sol:m1-2018s-q5}
12张的总子集数为$\binom{120}{12}$. 有利集合须分别从三种颜色选出指定数量, 有$\binom{30}{3}\binom{40}{4}\binom{50}{5}$种. 三部分颜色不同, 各自选择可唯一重组, 没有重复计数. 所求概率为
\[
\frac{\binom{30}{3}\binom{40}{4}\binom{50}{5}}{\binom{120}{12}}.
\]
取样无放回, 用各颜色的组合数比值计数, 不按12次独立取样处理.
\end{solution}

\begin{exercise}\label{ex:m1-2018s-q6}
\textbf{2018年春季, 原题6 (15分).} 10个人将帽子投入箱中随机混合后每人拿一顶. 若第$i$个人拿到自己的帽子, 令$X_i=1$, 否则$X_i=0$. 令$X=\sum_{i=1}^{10}X_i$为归位总人数. 计算:
\begin{enumerate}\item[(a)]\phantomsection\label{ex:m1-2018s-q6-a} $\E[X]$.\item[(b)]\phantomsection\label{ex:m1-2018s-q6-b} $\E[X_3X_7]$.\item[(c)]\phantomsection\label{ex:m1-2018s-q6-c} $\E[X_1^2+X_2^2+X_3^2]$.\end{enumerate}
\end{exercise}
\begin{solution}
\textbf{(a)}\phantomsection\label{sol:m1-2018s-q6-a} 每人归位概率为$1/10$, 故$\E[X]=\sum_{i=1}^{10}1/10=1$.

\textbf{(b)}\phantomsection\label{sol:m1-2018s-q6-b} 乘积$X_3X_7$是第3人与第7人都归位的指标. 固定这两人的帽子后, 其余8人的帽子可任意排列, 有$8!$种. 因而$\E[X_3X_7]=8!/10!=1/90$. 两个指标的边缘均值虽均为$1/10$, 它们不独立, 不能相乘成$1/100$.

\textbf{(c)}\phantomsection\label{sol:m1-2018s-q6-c} 指标只取0或1, 所以$X_i^2=X_i$. 期望线性给出$\E[X_1^2+X_2^2+X_3^2]=1/10+1/10+1/10=3/10$. 平方指标的计算不涉及两人同时归位的概率.
\end{solution}

\begin{exercise}\label{ex:m1-2018s-q7}
\textbf{2018年春季, 原题7 (15分).} Bob在机场小卖部想买一瓶没有标价的水. 他口渴却不好意思询价, 只知道美元价格$X$为3至7之间的整数, 并把$\{3,4,5,6,7\}$五个值各赋予$1/5$的概率. 按Bob的概率模型计算:
\begin{enumerate}\item[(a)]\phantomsection\label{ex:m1-2018s-q7-a} $\E[X]$.\item[(b)]\phantomsection\label{ex:m1-2018s-q7-b} $\Var[X]$.\item[(c)]\phantomsection\label{ex:m1-2018s-q7-c} $\Var[1.05X]$, 即假设机场所在州销售税为5\%时含税价格的方差.\end{enumerate}
\end{exercise}
\begin{solution}
\textbf{(a)}\phantomsection\label{sol:m1-2018s-q7-a} $\E[X]=(3+4+5+6+7)/5=5$.

\textbf{(b)}\phantomsection\label{sol:m1-2018s-q7-b} 将价格减去均值5, 偏差为$-2,-1,0,1,2$, 各概率$1/5$. 因而$\Var(X)=(4+1+0+1+4)/5=2$.

\textbf{(c)}\phantomsection\label{sol:m1-2018s-q7-c} 对常数$a$, $aX$相对其均值$a\E[X]$的偏差为$a(X-\E[X])$, 所以方差乘$a^2$. 取$a=1.05=21/20$, 得$\Var(1.05X)=(21/20)^2\cdot2=441/200$. 题设税率是原题的假定参数, 不替换成今日某州的税率.
\end{solution}

\originalsection{2019年春季}
% SOURCE 2019-spring-midterm1
\noindent 官方来源: \href{https://ocw.mit.edu/courses/18-600-probability-and-random-variables-fall-2019/1ef24cd4a1d900eabbc7e354063beebe_MIT18_600F19_mid1_2019.pdf}{原试卷PDF}; \href{https://ocw.mit.edu/courses/18-600-probability-and-random-variables-fall-2019/043a5f241c41a11ebbae7a84acfa94d2_MIT18_600F19_mid1_2019_soln.pdf}{官方解答PDF}.
\begin{exercise}\label{ex:m1-2019s-q1}
\textbf{2019年春季, 原题1 (20分).} 一个小镇有2000名居民, 唯一一家电影院正在放映一部冷门电影. 每位居民独立地以$1/1000$的概率去看这部电影. 设$X$为观看人数.
\begin{enumerate}\item[(a)]\phantomsection\label{ex:m1-2019s-q1-a} 精确计算$\E[X]$.\item[(b)]\phantomsection\label{ex:m1-2019s-q1-b} 精确计算$\Var[X]$.\item[(c)]\phantomsection\label{ex:m1-2019s-q1-c} 精确计算$\E[X^2]$.\item[(d)]\phantomsection\label{ex:m1-2019s-q1-d} 用泊松随机变量近似计算$\Prob(X=4)$.\end{enumerate}
\end{exercise}
\begin{solution}
\textbf{(a)}\phantomsection\label{sol:m1-2019s-q1-a} 观看人数是2000个独立观看指标之和, $X\sim\Bin(2000,1/1000)$, 故$\E[X]=2$.

\textbf{(b)}\phantomsection\label{sol:m1-2019s-q1-b} 每个指标方差为$(1/1000)(999/1000)$, 故$\Var(X)=2(999/1000)=999/500=1.998$.

\textbf{(c)}\phantomsection\label{sol:m1-2019s-q1-c} 方差恒等式给出$\E[X^2]=\Var(X)+\E[X]^2=1.998+4=5.998=2999/500$. 也可从指标$I_i$展开:
\[
\E[X^2]=\sum_{i=1}^{2000}\sum_{j=1}^{2000}\E[I_iI_j]
=2000p+(2000^2-2000)p^2,\qquad p=\frac1{1000}.
\]
这里$i=j$的2000项用$I_i^2=I_i$贡献$p$, $i\ne j$的项用独立性贡献$p^2$, 因而同样得到$2+3.998$. 官方解答最后双重和的两个下标原印都为$i=1$, 第二个应为$j=1$; 上式将不同指标及对角项明确分开.

\textbf{(d)}\phantomsection\label{sol:m1-2019s-q1-d} 用$\Pois(2)$近似, 得$\Prob(X=4)\approx\ee^{-2}2^4/4!=2\ee^{-2}/3$. (a)--(c)的有限小数是精确值, 此问才使用分布近似.
\end{solution}

\begin{exercise}\label{ex:m1-2019s-q2}
\textbf{2019年春季, 原题2 (10分).} $X$为参数2的泊松随机变量, $Y$为参数3的泊松随机变量.
\begin{enumerate}\item[(a)]\phantomsection\label{ex:m1-2019s-q2-a} 计算$\E[3X+4Y+5]$.\item[(b)]\phantomsection\label{ex:m1-2019s-q2-b} 计算$\Var(5X+7)$.\end{enumerate}
\end{exercise}
\begin{solution}
\textbf{(a)}\phantomsection\label{sol:m1-2019s-q2-a} 泊松分布的期望等于参数, $\E[X]=2$, $\E[Y]=3$. 期望线性给出$3\cdot2+4\cdot3+5=23$. 原题没有假设$X,Y$独立, 此处也不需要它.

\textbf{(b)}\phantomsection\label{sol:m1-2019s-q2-b} 加常数7不改变偏差, 乘5使平方偏差乘25. 泊松分布的方差等于参数, 所以$\Var(5X+7)=25\Var(X)=25\cdot2=50$. 本问仅含$X$, 不涉及$X,Y$的协方差.
\end{solution}

\begin{exercise}\label{ex:m1-2019s-q3}
\textbf{2019年春季, 原题3 (20分).} Alice、Bob、Carol、Dave、Eve和Frank聚在一起吃披萨并玩《龙与地下城》. 他们买了两个大披萨, 每个切12片, 共24片.
\begin{enumerate}
\item[(a)]\phantomsection\label{ex:m1-2019s-q3-a} 把24片不可区分的披萨分给6个人, 有多少种方式? 即有多少个非负整数序列$a_1,\ldots,a_6$满足$\sum_{i=1}^6a_i=24$?
\item[(b)]\phantomsection\label{ex:m1-2019s-q3-b} Eve提议只考虑每人至少1片的分法以求公平. 有多少个严格正整数序列$a_1,\ldots,a_6$满足$\sum_{i=1}^6a_i=24$?
\item[(c)]\phantomsection\label{ex:m1-2019s-q3-c} 6个玩家各独立掷一个公平20面骰, 点数为$\{1,2,\ldots,20\}$. 求六个点数之和恰为24的概率.
\end{enumerate}
\end{exercise}
\begin{solution}
\textbf{(a)}\phantomsection\label{sol:m1-2019s-q3-a} 24颗星和5条分隔线编码6个人的片数, 零片对应空段. 共29个位置, 分配数为$\binom{29}{5}$.

\textbf{(b)}\phantomsection\label{sol:m1-2019s-q3-b} 先每人发1片, 剩下18片任意分配. 对$b_i=a_i-1\ge0$, 有$\sum b_i=18$, 故星与杠给出$\binom{18+6-1}{6-1}=\binom{23}{5}$.

\textbf{(c)}\phantomsection\label{sol:m1-2019s-q3-c} 骰子点数须在1至20. (b)中任一正整数解的单个分量至多$24-5=19$, 因为另5个分量各至少1. 因而20的上界在总和24时自动满足, (b)的全部解正好是此问有利点数序列. 独立公平骰共有$20^6$个等可能序列, 所以概率为$\binom{23}{5}/20^6$. 若总和更大而允许单个分量超过20, 就不能直接沿用无上界的星与杠计数.
\end{solution}

\begin{exercise}\label{ex:m1-2019s-q4}
\textbf{2019年春季, 原题4 (20分).} 一个有15名成员的无伴奏合唱团, 其中8名女性、7名男性, 举办节日礼物交换. 每人把姓名写在纸上投入碗中, 再随机把15张姓名纸分给15个人, 全部$15!$种安排等可能. 每人给自己抽到姓名的那个人买礼物.
\begin{enumerate}\item[(a)]\phantomsection\label{ex:m1-2019s-q4-a} 计算被分配给自己买礼物的人数的期望.\item[(b)]\phantomsection\label{ex:m1-2019s-q4-b} 计算被分配给女性送礼的男性人数的期望.\item[(c)]\phantomsection\label{ex:m1-2019s-q4-c} 计算每一位男性都被分配给女性送礼的概率.\item[(d)]\phantomsection\label{ex:m1-2019s-q4-d} 计算每个人都属于长度3循环的概率, 即每组中的$A$给$B$, $B$给$C$, $C$给$A$.\end{enumerate}
\end{exercise}
\begin{solution}
\textbf{(a)}\phantomsection\label{sol:m1-2019s-q4-a} 每人抽到自己名字的概率为$1/15$. 将15个自送礼指标的均值加起来得$15/15=1$.

\textbf{(b)}\phantomsection\label{sol:m1-2019s-q4-b} 任一男性抽到8个女性名字中的一个的概率为$8/15$. 对7位男性的指标求和, 期望为$7(8/15)=56/15$; 不需要男性之间的指标独立.

\textbf{(c)}\phantomsection\label{sol:m1-2019s-q4-c} 依次让7位男性抽签, 第一位有8个女性名字可抽, 第二位剩7个, 依此第7位剩2个. 因而所求为
\[
\frac8{15}\frac7{14}\frac6{13}\frac5{12}\frac4{11}\frac3{10}\frac2{9}
=\frac{(8!/1!)}{(15!/8!)}=\frac{(8!)^2}{15!}.
\]
男性各抽到不同女性后, 剩余8个名字如何分给女性都允许, 没有附加男性必须收到女性礼物的条件.

\textbf{(d)}\phantomsection\label{sol:m1-2019s-q4-d} 15个人须分成5个无序三人组. 先排15人再每3人一组, 除去每组内部$3!$种顺序及5组的$5!$种顺序, 组划分数为$15!/((3!)^5 5!)$. 每三人组有两个不同的有向三循环, 所以有利排列数为$2^5\,15!/((3!)^5 5!)$. 除以$15!$得概率$2^5/((3!)^5 5!)=1/(3^5 5!)$. 组内的方向因子2不能省略, 也不能把三循环按三个可能起点再次计数.
\end{solution}

\begin{exercise}\label{ex:m1-2019s-q5}
\textbf{2019年春季, 原题5 (10分).} 标准52张牌中, 方块、红桃、梅花、黑桃各13张. 随机分成4手桥牌, 每手13张, 所有分配等可能. 求每一手的牌都只属于一种花色的概率, 即一手全红桃, 一手全梅花, 依此类推.
\end{exercise}
\begin{solution}
\phantomsection\label{sol:m1-2019s-q5}
将4手视为给4位不同玩家, 全部分配数为$52!/(13!)^4$. 有利分配须将4种花色一一指定给4位玩家, 共$4!$种; 各指定花色的全部13张必须进入对应手牌, 不再有其他选择. 所以概率为$4!(13!)^4/52!$. 若把4手视为无序分组, 总数与有利数同时除以$4!$, 概率不变.
\end{solution}

\begin{exercise}\label{ex:m1-2019s-q6}
\textbf{2019年春季, 原题6 (20分).} Alicia写历史课论文. 每次论文以0.7的概率很出色, 以0.3的概率平庸. 教授读到出色论文时以0.9的概率给A, 读到平庸论文时以0.3的概率给A. 设$B$为论文出色, $A$为获得A成绩, 即$\Prob(B)=0.7$, $\Prob(A\mid B)=0.9$, $\Prob(A\mid B^c)=0.3$.
\begin{enumerate}\item[(a)]\phantomsection\label{ex:m1-2019s-q6-a} 计算$\Prob(A)$, 即总体获得A的概率.\item[(b)]\phantomsection\label{ex:m1-2019s-q6-b} 计算$\Prob(B\mid A)$, 即获得A后论文出色的概率.\item[(c)]\phantomsection\label{ex:m1-2019s-q6-c} 计算$\Prob(B\mid A^c)$, 即未获得A后论文出色的概率.\end{enumerate}
\end{exercise}
\begin{solution}
\textbf{(a)}\phantomsection\label{sol:m1-2019s-q6-a} 按出色与平庸分解, $\Prob(A)=0.7\cdot0.9+0.3\cdot0.3=0.63+0.09=0.72$.

\textbf{(b)}\phantomsection\label{sol:m1-2019s-q6-b} $\Prob(BA)=0.63$, 所以$\Prob(B\mid A)=0.63/0.72=7/8$.

\textbf{(c)}\phantomsection\label{sol:m1-2019s-q6-c} 出色且未获A的概率为$0.7(1-0.9)=0.07$, 未获A的总概率为$1-0.72=0.28$. 因而$\Prob(B\mid A^c)=0.07/0.28=1/4$. 给出的出色率与两个评分条件概率均保留, 不把评分事件假定为与论文质量独立.
\end{solution}

\originalsection{2019年秋季}
% SOURCE 2019-fall-midterm1
\noindent 官方来源: \href{https://ocw.mit.edu/courses/18-600-probability-and-random-variables-fall-2019/dea9112a7fe28281e2c236ab6334f7ad_MIT18_600F19_midterm1.pdf}{原试卷PDF}; \href{https://ocw.mit.edu/courses/18-600-probability-and-random-variables-fall-2019/62e1b919e9f7bf208a943bde0d0dc607_MIT18_600F19_midterm1_soln.pdf}{官方解答PDF}.
\begin{exercise}\label{ex:m1-2019f-q1}
\textbf{2019年秋季, 原题1 (15分).} 超级喷发是产生超过1000立方千米沉积物的火山喷发, 灰尘还可能使全球气候改变数年. 假设每一年独立地以$1/25000$的概率, 世界上某处恰发生一次超级喷发. 为简化问题, 假设同一年发生超过一次超级喷发的概率为零.
\begin{enumerate}\item[(a)]\phantomsection\label{ex:m1-2019f-q1-a} 计算未来100000年内超级喷发次数的期望.\item[(b)]\phantomsection\label{ex:m1-2019f-q1-b} 用泊松近似估计未来100000年内恰有3次超级喷发的概率.\item[(c)]\phantomsection\label{ex:m1-2019f-q1-c} 用泊松近似估计未来100年内至少出现一次超级喷发的概率.\end{enumerate}
\end{exercise}
\begin{solution}
\textbf{(a)}\phantomsection\label{sol:m1-2019f-q1-a} 年喷发数为取值0或1的独立指标. 100000年的总数服从$\Bin(100000,1/25000)$. 其期望为$100000/25000=4$.

\textbf{(b)}\phantomsection\label{sol:m1-2019f-q1-b} 用均值4的$\Pois(4)$近似, 恰3次的概率为$\ee^{-4}4^3/3!=32\ee^{-4}/3$.

\textbf{(c)}\phantomsection\label{sol:m1-2019f-q1-c} 此问时间缩为100年, 所以参数须重新计算为$100/25000=1/250$. 用$\Pois(1/250)$近似, 至少一次的概率为$1-\ee^{-1/250}$. 这是从零次事件取补集所得, 不能沿用(b)的参数4. 若进一步作一阶近似, $1-\ee^{-1/250}\approx1/250$, 但题目所要求的泊松近似保留指数形式即可.
\end{solution}

\begin{exercise}\label{ex:m1-2019f-q2}
\textbf{2019年秋季, 原题2 (15分).} 两队参加没有加时或点球大战的足球赛. 第一队进球数$X$为参数$\lambda_X=1$的泊松随机变量, 第二队进球数$Y$为参数$\lambda_Y=2$的独立泊松随机变量.
\begin{enumerate}\item[(a)]\phantomsection\label{ex:m1-2019f-q2-a} 计算平局概率$\Prob(X=Y)$, 可以保留无限和.\item[(b)]\phantomsection\label{ex:m1-2019f-q2-b} 计算较弱的第一队获胜的概率$\Prob(X>Y)$, 可以保留双重无限和.\item[(c)]\phantomsection\label{ex:m1-2019f-q2-c} 计算整场恰有2球的概率$\Prob(X+Y=2)$.\end{enumerate}
\end{exercise}
\begin{solution}
\textbf{(a)}\phantomsection\label{sol:m1-2019f-q2-a} 平局由$(X,Y)=(k,k)$的互斥事件组成, $k\ge0$. 独立性给出
\[
\Prob(X=Y)=\sum_{k=0}^{\infty}
\left(\ee^{-1}\frac{1^k}{k!}\right)
\left(\ee^{-2}\frac{2^k}{k!}\right)
=\ee^{-3}\sum_{k=0}^{\infty}\frac{2^k}{(k!)^2}.
\]
级数每项非负, 比值为$2/(k+1)^2\to0$, 因而收敛, 也可直接由它是一系列互斥事件概率之和看出不超过1.

\textbf{(b)}\phantomsection\label{sol:m1-2019f-q2-b} 固定第一队进球$j\ge1$, 第二队须进$k=0,\ldots,j-1$球. 把这些互斥比分求和得
\[
\Prob(X>Y)=\sum_{j=1}^{\infty}\sum_{k=0}^{j-1}
\left(\ee^{-1}\frac1{j!}\right)
\left(\ee^{-2}\frac{2^k}{k!}\right)
=\ee^{-3}\sum_{j=1}^{\infty}\frac1{j!}\sum_{k=0}^{j-1}\frac{2^k}{k!}.
\]
内和有界于$\ee^2$, 外和有界于$\ee-1$, 所以式子收敛. 上界$k=j-1$体现严格获胜, 没有把平局算进来.

\textbf{(c)}\phantomsection\label{sol:m1-2019f-q2-c} 仅比分$(0,2),(1,1),(2,0)$有利. 由独立性各概率依次为$2\ee^{-3}$、$2\ee^{-3}$、$\ee^{-3}/2$, 加起来为$9/(2\ee^3)$. 也可验证独立泊松和$X+Y\sim\Pois(3)$, 其恰2球概率为$\ee^{-3}3^2/2!$, 与直接计数一致.
\end{solution}

\begin{exercise}\label{ex:m1-2019f-q3}
\textbf{2019年秋季, 原题3 (20分).} 化学课有14名学生, 教授给每人指定一位伙伴, 形成7个每组2人的无序配对.
\begin{enumerate}
\item[(a)]\phantomsection\label{ex:m1-2019f-q3-a} 共有多少种配对方式?
\item[(b)]\phantomsection\label{ex:m1-2019f-q3-b} 两名学生是好友, 希望成为伙伴. 若教授在所有配对方式中均匀随机选择, 他们配为伙伴的概率是多少?
\item[(c)]\phantomsection\label{ex:m1-2019f-q3-c} 假设7名学生为男性、7名为女性. 令$N$为一男一女的配对数, 计算$\E[N]$. 如有帮助, 可令$N_i$为第$i$名女性有男性伙伴时取1, 否则取0的随机变量.
\item[(d)]\phantomsection\label{ex:m1-2019f-q3-d} 计算$\E[N^2]$.
\end{enumerate}
\end{exercise}
\begin{solution}
\textbf{(a)}\phantomsection\label{sol:m1-2019f-q3-a} 将14人排列并相邻配对. 每个配对方案由于7对内部两人的次序和7对之间的次序, 被计数$2^7 7!$次. 因而方案数为$14!/(2^7 7!)=13\cdot11\cdot9\cdot7\cdot5\cdot3\cdot1$.

\textbf{(b)}\phantomsection\label{sol:m1-2019f-q3-b} 固定其中一位好友, 其伙伴在其余13人中等可能, 所以概率为$1/13$. 也可固定好友配对后数剩余12人的方案$12!/(2^6 6!)$, 除以(a)得到同一值.

\textbf{(c)}\phantomsection\label{sol:m1-2019f-q3-c} 每个异性配对恰含1名女性, 所以$N=\sum_{i=1}^7N_i$, 不会把异性对数两次. 一位女性在其他13人中有7个男性可作伙伴, 故$\E[N_i]=7/13$, $\E[N]=7(7/13)=49/13$.

\textbf{(d)}\phantomsection\label{sol:m1-2019f-q3-d} 平方和展开为$N^2=\sum_iN_i^2+\sum_{i\ne j}N_iN_j$. 因为$N_i^2=N_i$, 7个对角项总贡献$49/13$. 对不同女性$i,j$, 先令$i$与某位男性配对, 概率为$7/13$; 移去这两人后, 剩12人中有6男6女, 女性$j$在其余11人中有6名男性可选. 条件于任何一个固定的$i$的男性伙伴, 剩余配对仍均匀, 所以
\[
\E[N_iN_j]=\frac7{13}\frac6{11},\qquad i\ne j.
\]
有$7\cdot6=42$个有序不同下标项, 因而
\[
\E[N^2]=\frac{49}{13}+42\frac7{13}\frac6{11}
=\frac{2303}{143}.
\]
不同女性的伙伴指标不独立, 第二因子的分母是11而不是13, 这是二阶矩与一阶矩计算的区别.
\end{solution}

\begin{exercise}\label{ex:m1-2019f-q4}
\textbf{2019年秋季, 原题4 (20分).} 计算:
\begin{enumerate}
\item[(a)]\phantomsection\label{ex:m1-2019f-q4-a} $\displaystyle\lim_{n\to\infty}\left(1-\frac1{4n}\right)^n$.
\item[(b)]\phantomsection\label{ex:m1-2019f-q4-b} $\displaystyle\sum_{k=0}^{9}2^k8^{9-k}\binom9k$.
\item[(c)]\phantomsection\label{ex:m1-2019f-q4-c} $\displaystyle\sum_{k=0}^{\infty}\frac1{2^k k!}$.
\item[(d)]\phantomsection\label{ex:m1-2019f-q4-d} $\displaystyle\sum_{k=0}^{\infty}\left(\frac56\right)^{k-1}\left(\frac16\right)k$.
\end{enumerate}
\end{exercise}
\begin{solution}
\textbf{(a)}\phantomsection\label{sol:m1-2019f-q4-a} 对每个正整数$n$, 底数为正, 可取对数. 设$t=-1/(4n)$, 则$n\ln(1-1/(4n))=-(1/4)\ln(1+t)/t\to-1/4$. 指数函数连续, 原极限为$\ee^{-1/4}$.

\textbf{(b)}\phantomsection\label{sol:m1-2019f-q4-b} 二项式定理的第$k$项为$\binom9k2^k8^{9-k}$, 因而全和就是$(2+8)^9=10^9$.

\textbf{(c)}\phantomsection\label{sol:m1-2019f-q4-c} 将$1/2^k$写为$(1/2)^k$, 识别指数级数: $\sum_{k\ge0}(1/2)^k/k!=\ee^{1/2}$. 比值$1/[2(k+1)]\to0$保证级数收敛.

\textbf{(d)}\phantomsection\label{sol:m1-2019f-q4-d} 原PDF的$k$是与$1/6$相乘, 并不是$1/6$的指数. 原下限确为$k=0$, 该项等于$(5/6)^{-1}(1/6)\cdot0=0$, 所以也可从$k=1$起求和. 令$q=5/6$, $p=1/6$. 几何等待时间$G$满足$\Prob(G=k)=q^{k-1}p$, $k\ge1$, 因而此和为$\E[G]=1/p=6$. 若直接用级数计算, 在$|q|<1$内部对$\sum_{k\ge0}q^k=1/(1-q)$求导, 得
\[
\sum_{k=1}^{\infty}kq^{k-1}=\frac1{(1-q)^2},\qquad
\sum_{k=0}^{\infty}q^{k-1}pk=\frac{p}{(1-q)^2}=6.
\]
两种推导核对了乘法、求和起点及收敛条件.
\end{solution}

\begin{exercise}\label{ex:m1-2019f-q5}
\textbf{2019年秋季, 原题5 (15分).} 烘焙比赛有10名参赛者. 评委对大多数烘焙材料过敏, 因而不品尝食物, 改为随机选择3位获胜者, 全部三人子集等可能. Alice、Bob和Carol三人是好友, 希望都获胜, 又觉得只有两人获胜而第三人没获胜会很尴尬. 设$A$为三人全部获胜, $B$为三人中至少两人获胜.
\begin{enumerate}\item[(a)]\phantomsection\label{ex:m1-2019f-q5-a} 计算$\Prob(A)$.\item[(b)]\phantomsection\label{ex:m1-2019f-q5-b} 计算$\Prob(B)$.\item[(c)]\phantomsection\label{ex:m1-2019f-q5-c} 计算$\Prob(A\mid B)$.\end{enumerate}
\end{exercise}
\begin{solution}
\textbf{(a)}\phantomsection\label{sol:m1-2019f-q5-a} 全部赢家集合有$\binom{10}{3}=120$种, 三位好友全胜仅有$\{\text{Alice,Bob,Carol}\}$这一种. 故$\Prob(A)=1/120$.

\textbf{(b)}\phantomsection\label{sol:m1-2019f-q5-b} 三人全胜有1种; 恰2人胜, 选出其中2人有$\binom32=3$种, 再从其余7人选第3位赢家, 有7种. 两种情况互斥, 因而$\Prob(B)=(1+3\cdot7)/120=22/120=11/60$.

\textbf{(c)}\phantomsection\label{sol:m1-2019f-q5-c} $A\subseteq B$, 所以$\Prob(AB)=\Prob(A)$, 条件概率为$(1/120)/(22/120)=1/22$. 条件是“至少两位好友获胜”, 不是已经指定哪两位获胜; (b)的22种结果必须全部留在分母中.
\end{solution}

\begin{exercise}\label{ex:m1-2019f-q6}
\textbf{2019年秋季, 原题6 (15分).} Janet觉得自己或许发烧, 或许只是头痛或感冒, 尚不确定. 她将用一个显示华氏温度、四舍五入到最接近整数的数字体温计量体温, 并认为会以各$1/5$的概率看到$\{98,99,100,101,102\}$中的一个值. 设$X$为实际显示的数字.
\begin{enumerate}
\item[(a)]\phantomsection\label{ex:m1-2019f-q6-a} 计算$\Var(X)$, 化简为一个明确的有理数精确值.
\item[(b)]\phantomsection\label{ex:m1-2019f-q6-b} 用(a)的答案计算$\E[X^2]$.
\item[(c)]\phantomsection\label{ex:m1-2019f-q6-c} 量过体温后, 无论显示什么, Janet都将掷一个三面骰, 其值$\{0,1,2\}$等可能. 如果点数为$N$, 她将服用$N$片布洛芬. 计算$\E[2N^3+3X^2]$.
\end{enumerate}
\end{exercise}
\begin{solution}
\textbf{(a)}\phantomsection\label{sol:m1-2019f-q6-a} 五个温度以100为中心对称, $\E[X]=(98+99+100+101+102)/5=100$. 平方偏差为$4,1,0,1,4$, 所以$\Var(X)=(4+1+0+1+4)/5=2$.

\textbf{(b)}\phantomsection\label{sol:m1-2019f-q6-b} $\Var(X)=\E[X^2]-\E[X]^2$, 因而$\E[X^2]=2+100^2=10002$.

\textbf{(c)}\phantomsection\label{sol:m1-2019f-q6-c} 对三面骰取平均, $\E[N^3]=(0^3+1^3+2^3)/3=3$. 期望线性给出$\E[2N^3+3X^2]=2\cdot3+3\cdot10002=30012$. 此处没有混合乘积, 只需给定两个边缘分布, 不需额外添加$N$与$X$独立的假设; 温度与片数仍是原试题的概率模型.
\end{solution}
